\documentclass[10pt]{article}
\usepackage[margin=1in]{geometry}
\usepackage[T1]{fontenc}
\usepackage[utf8]{inputenc}
\usepackage{hyperref}

\title{Jurassic Design Motifs:\\
Extended Abstract}
\author{Bitcoin Consensus Observatory (Jurassic Bitcoin)}
\date{\today}

\begin{document}
\maketitle

\begin{abstract}
Bitcoin's pre-SegWit consensus history contains behaviors that are no longer
desirable transaction forms but remain informative as design motifs. We focus
on three reusable motifs: transcript multiplicity, identifier bifurcation, and
carrier camouflage. Our contribution is an observer pipeline that extracts
historical Bitcoin windows, materializes seam-family fixtures, classifies them
with a policy-envelope map, and then lifts the resulting abstractions into live
TradeLayer Litecoin prototypes with BitVM-style dispute flows
\cite{linus2023bitvm,linus2026binohash,yang2021fluffy,kim2025fork}. The point
is not to revive old bugs. It is to translate fossilized Bitcoin behavior into
measured overlay search surfaces.
\end{abstract}

\section*{Motifs and Method}

We use \emph{Jurassic Bitcoin} as shorthand for the pre-SegWit historical era
in which Bitcoin still exposed many laissez-faire technical fingerprints of its
early libertarian design culture. That era is over, but some of those old
fingerprints survive as reusable protocol shapes once they are abstracted away
from the original transaction forms that carried them. In that sense, the
observer pipeline is the mosquito in amber: it preserves small but
information-rich traces of a vanished protocol ecology, then lets us recover
isomorphic mechanics without reviving the original environment itself.

Transcript multiplicity captures cases where one spend family admits multiple
effective signing transcripts. Identifier bifurcation captures cases where a
public handle moves while the committed semantic core stays fixed. Carrier
camouflage captures the idea that overlay commitments can hide inside
ordinary-looking transaction topology rather than only in explicit exotic
carriers. Historically, these families are visible in legacy sighash seams,
dummy-element txid variation, and 2013 payout-shaped transaction bursts, while
modern analogues are emerging in BitVM-adjacent introspection work and other
transaction-shape-aware constructions \cite{linus2026binohash,andrychowicz2013malleability,decker2014malleability,ghesmati2022payjoin}.

Methodologically, the repository acts as an executable observer. It extracts
activation-era windows, constructs deterministic fixtures, runs replay and
classification passes, and maps each surface into one of three buckets:
replay-only fossil, ordinary historical cover, or clean modern isomorphism.
That observer stance is closer to differential blockchain testing than to a
paper-only taxonomy \cite{yang2021fluffy,kim2025fork}.

\section*{Prototype Results}

The strongest current live evidence comes from Litecoin/TradeLayer procedural
token prototypes. We have live proofs for transcript alias relays, identifier
namespace relays, and a fused router-dispute flow that combines transcript and
namespace preludes with BitVM-style cache, challenge, resolve, and payout
branches. One branch releases funds and another remains blocked, showing that
the motifs are operational design patterns rather than descriptive labels
\cite{linus2023bitvm}.

The repository now also includes three non-BitVM application meshes on LTCTEST:
an oracle sidecar mesh, a watchtower beacon mesh, and a statechain handoff
mesh. In each case, one semantic state is expressed through multiple relay
transcripts, rotated public handles, and ordinary-looking carrier hints. That
extension matters because it shows the three motifs are not only useful for
BitVM-flavored dispute graphs; they also transfer to oracle publication,
monitoring, and checkpoint systems.

The Bitcoin-side oracle-sidecar replay bench is partially complete: the
historical carrier families are extracted and synthetic replay fixtures exist,
but live Bitcoin materialization still needs funded prevouts. That makes the
paper's main claim deliberately modest. Jurassic Bitcoin is best understood as a
source of reusable shapes, not as a menu of deployable legacy mechanisms.

\section*{Future Work}

The immediate next steps are to materialize funded Bitcoin-side replay inputs,
turn the Litecoin graphs into reusable procedural-token templates, and encourage
other developers to attempt Binohash-style adaptations of transcript aliasing,
namespace indirection, and ordinary-topology commitment cover in their own DLC,
BitVM, and contract-layer constructions \cite{linus2026binohash}. The intended
audience is broad: nonprofit-supported researchers, corporate protocol teams,
independent maintainers, and new contributors can all treat these motifs as a
small experimental toolkit rather than a fixed agenda, and try them in their
own systems if they prove directionally useful.

\bibliographystyle{plain}
\bibliography{jurassic_design_motifs}

\end{document}
